\documentclass[10pt,a4paper]{amsart}
\usepackage{amsmath,amssymb,amsthm,graphicx}
\usepackage{verbatim}
\newenvironment{reference}{\comment}{\endcomment}
%\newenvironment{reference}{}{}
\newenvironment{slogan}{\comment}{\endcomment}
\newenvironment{history}{\comment}{\endcomment}

% For commutative diagrams we use Xy-pic
\usepackage[all]{xy}

% We use 2cell for 2-commutative diagrams.
\xyoption{2cell}
% Theorem environments.
%
\theoremstyle{plain}
\newtheorem{theorem}[subsection]{Theorem}
\newtheorem{proposition}[subsection]{Proposition}
\newtheorem{lemma}[subsection]{Lemma}

\theoremstyle{definition}
\newtheorem{definition}[subsection]{Definition}
\newtheorem{example}[subsection]{Example}
\newtheorem{exercise}[subsection]{Exercise}
\newtheorem{situation}[subsection]{Situation}

\theoremstyle{remark}
\newtheorem{remark}[subsection]{Remark}
\newtheorem{remarks}[subsection]{Remarks}

\numberwithin{equation}{subsection}

% Macros
%
\def\lim{\mathop{\mathrm{lim}}\nolimits}
\def\colim{\mathop{\mathrm{colim}}\nolimits}
\def\Spec{\mathop{\mathrm{Spec}}}
\def\Hom{\mathop{\mathrm{Hom}}\nolimits}
\def\Ext{\mathop{\mathrm{Ext}}\nolimits}
\def\SheafHom{\mathop{\mathcal{H}\!\mathit{om}}\nolimits}
\def\SheafExt{\mathop{\mathcal{E}\!\mathit{xt}}\nolimits}
\def\Sch{\mathit{Sch}}
\def\Mor{\mathop{\mathrm{Mor}}\nolimits}
\def\Ob{\mathop{\mathrm{Ob}}\nolimits}
\def\Sh{\mathop{\mathit{Sh}}\nolimits}
\def\NL{\mathop{N\!L}\nolimits}
\def\CH{\mathop{\mathrm{CH}}\nolimits}
\def\proetale{{pro\text{-}\acute{e}tale}}
\def\etale{{\acute{e}tale}}
\def\QCoh{\mathit{QCoh}}
\def\Ker{\mathop{\mathrm{Ker}}}
\def\Im{\mathop{\mathrm{Im}}}
\def\Coker{\mathop{\mathrm{Coker}}}
\def\Coim{\mathop{\mathrm{Coim}}}

% Boxtimes
%
\DeclareMathSymbol{\boxtimes}{\mathbin}{AMSa}{"02}

%
% Macros for moduli stacks/spaces
%
\def\QCohstack{\mathcal{QC}\!\mathit{oh}}
\def\Cohstack{\mathcal{C}\!\mathit{oh}}
\def\Spacesstack{\mathcal{S}\!\mathit{paces}}
\def\Quotfunctor{\mathrm{Quot}}
\def\Hilbfunctor{\mathrm{Hilb}}
\def\Curvesstack{\mathcal{C}\!\mathit{urves}}
\def\Polarizedstack{\mathcal{P}\!\mathit{olarized}}
\def\Complexesstack{\mathcal{C}\!\mathit{omplexes}}
% \Pic is the operator that assigns to X its picard group, usage \Pic(X)
% \Picardstack_{X/B} denotes the Picard stack of X over B
% \Picardfunctor_{X/B} denotes the Picard functor of X over B
\def\Pic{\mathop{\mathrm{Pic}}\nolimits}
\def\Picardstack{\mathcal{P}\!\mathit{ic}}
\def\Picardfunctor{\mathrm{Pic}}
\def\Deformationcategory{\mathcal{D}\!\mathit{ef}}
\usepackage{cleveref}
\crefname{originallemma}{original lemma}{original lemmas}
\Crefname{originallemma}{Original lemma}{Original lemmas}
\makeatletter
\newlabel{lemma-lisse-etale}{{0788}{source}{}{sourcepointer-1}{}}
\newlabel{lemma-lisse-etale@cref}{{[originallemma][0][]0788}{[1][0][]source}{}{}{}}
\makeatother
\crefname{originallemma}{original lemma}{original lemmas}
\Crefname{originallemma}{Original lemma}{Original lemmas}
\makeatletter
\newlabel{lemma-adjoint}{{0778}{source}{}{sourcepointer-2}{}}
\newlabel{lemma-adjoint@cref}{{[originallemma][0][]0778}{[1][0][]source}{}{}{}}
\makeatother
\crefname{originalreference}{original reference}{original references}
\Crefname{originalreference}{Original reference}{Original references}
\makeatletter
\newlabel{sites-modules-lemma-local-pullback}{{03DO}{source}{}{sourcepointer-3}{}}
\newlabel{sites-modules-lemma-local-pullback@cref}{{[originalreference][0][]03DO}{[1][0][]source}{}{}{}}
\makeatother
\crefname{originalreference}{original reference}{original references}
\Crefname{originalreference}{Original reference}{Original references}
\makeatletter
\newlabel{stacks-sheaves-lemma-characterize-quasi-coherent-bis}{{0EM8}{source}{}{sourcepointer-4}{}}
\newlabel{stacks-sheaves-lemma-characterize-quasi-coherent-bis@cref}{{[originalreference][0][]0EM8}{[1][0][]source}{}{}{}}
\makeatother
\crefname{originallemma}{original lemma}{original lemmas}
\Crefname{originallemma}{Original lemma}{Original lemmas}
\makeatletter
\newlabel{lemma-shriek-quasi-coherent}{{07B0}{source}{}{sourcepointer-5}{}}
\newlabel{lemma-shriek-quasi-coherent@cref}{{[originallemma][0][]07B0}{[1][0][]source}{}{}{}}
\makeatother
\crefname{originalreference}{original reference}{original references}
\Crefname{originalreference}{Original reference}{Original references}
\makeatletter
\newlabel{stacks-sheaves-section-quasi-coherent-presentation}{{06WS}{source}{}{sourcepointer-6}{}}
\newlabel{stacks-sheaves-section-quasi-coherent-presentation@cref}{{[originalreference][0][]06WS}{[1][0][]source}{}{}{}}
\makeatother
\crefname{originalreference}{original reference}{original references}
\Crefname{originalreference}{Original reference}{Original references}
\makeatletter
\newlabel{stacks-sheaves-section-quasi-coherent-algebraic-stacks}{{06WU}{source}{}{sourcepointer-7}{}}
\newlabel{stacks-sheaves-section-quasi-coherent-algebraic-stacks@cref}{{[originalreference][0][]06WU}{[1][0][]source}{}{}{}}
\makeatother
\crefname{originallemma}{original lemma}{original lemmas}
\Crefname{originallemma}{Original lemma}{Original lemmas}
\makeatletter
\newlabel{lemma-adjoint-kernel-parasitic}{{0779}{source}{}{sourcepointer-8}{}}
\newlabel{lemma-adjoint-kernel-parasitic@cref}{{[originallemma][0][]0779}{[1][0][]source}{}{}{}}
\makeatother
\crefname{originallemma}{original lemma}{original lemmas}
\Crefname{originallemma}{Original lemma}{Original lemmas}
\makeatletter
\newlabel{lemma-parasitic-in-terms-flat-fppf}{{07AR}{source}{}{sourcepointer-9}{}}
\newlabel{lemma-parasitic-in-terms-flat-fppf@cref}{{[originallemma][0][]07AR}{[1][0][]source}{}{}{}}
\makeatother
\title{g-inverse and g-shriek identify QCoh with qc modules on the lisse-etale and flat-fppf sites}
\author{The Stacks Project Authors}
\date{stacks-project:stacks-cohomology.tex:local-snapshot-20261012; GNU FDL 1.2 or later}
\begin{document}
\maketitle
\noindent\textit{Editorial note.} All delivered mathematical text is the author's own, copied verbatim from the pinned Stacks Project chapter file stacks-cohomology.tex; no statement or proof step was rewritten, reordered or removed.\par\medskip
\noindent\textit{Editorial note.} The counted span cites the lemmas recorded in assembly.external\_references with canonical tags as navigation pointers.\par\medskip
\noindent\textit{Editorial note.} Spot-check note: The isomorphism g\_!g\textasciicircum{}\{-1\}F -> F in (a) is concluded by 'the discussion in Sheaves on Stacks, Sections ... shows it is an isomorphism' - a section-discussion reference rather than a numbered lemma.\par\medskip
\noindent\textit{Allowed standard prerequisites (editorial).}\begin{itemize}
\item The lisse-etale and flat-fppf morphisms of topoi. Cohomology of Algebraic Stacks, Lemma stacks-cohomology-lemma-lisse-etale. g\textasciicircum{}* = g\textasciicircum{}\{-1\} and g\textasciicircum{}\{-1\}g\_! = id on sheaves.
\item Q-adjunction defect is parasitic. Cohomology of Algebraic Stacks, Lemma stacks-cohomology-lemma-adjoint-kernel-parasitic (2). Kernel/cokernel parasitic.
\item Parasitic detection on the sites. Cohomology of Algebraic Stacks, Lemma stacks-cohomology-lemma-parasitic-in-terms-flat-fppf. g\textasciicircum{}* of a parasitic module is zero.
\end{itemize}\medskip
\ifdefined\hypertarget\hypertarget{sourcepointer-1}{}\fi
\noindent\textit{Source reference.} Original lemma 0788: sources/stacks-chow/tags.txt, source lines 9228--9228. Cross-chapter standard prerequisite cited by the author inside the counted span; it is used as a pointer and its own proof is not part of this item.\par\medskip
\ifdefined\hypertarget\hypertarget{sourcepointer-2}{}\fi
\noindent\textit{Source reference.} Original lemma 0778: sources/stacks-chow/tags.txt, source lines 9193--9193. Cross-chapter standard prerequisite cited by the author inside the counted span; it is used as a pointer and its own proof is not part of this item.\par\medskip
\ifdefined\hypertarget\hypertarget{sourcepointer-3}{}\fi
\noindent\textit{Source reference.} Original reference 03DO: sources/stacks-chow/tags.txt, source lines 4427--4427. Cross-chapter standard prerequisite cited by the author inside the counted span; it is used as a pointer and its own proof is not part of this item.\par\medskip
\ifdefined\hypertarget\hypertarget{sourcepointer-4}{}\fi
\noindent\textit{Source reference.} Original reference 0EM8: sources/stacks-chow/tags.txt, source lines 18298--18298. Cross-chapter standard prerequisite cited by the author inside the counted span; it is used as a pointer and its own proof is not part of this item.\par\medskip
\ifdefined\hypertarget\hypertarget{sourcepointer-5}{}\fi
\noindent\textit{Source reference.} Original lemma 07B0: sources/stacks-chow/tags.txt, source lines 9325--9325. Cross-chapter standard prerequisite cited by the author inside the counted span; it is used as a pointer and its own proof is not part of this item.\par\medskip
\ifdefined\hypertarget\hypertarget{sourcepointer-6}{}\fi
\noindent\textit{Source reference.} Original reference 06WS: sources/stacks-chow/tags.txt, source lines 8826--8826. Cross-chapter standard prerequisite cited by the author inside the counted span; it is used as a pointer and its own proof is not part of this item.\par\medskip
\ifdefined\hypertarget\hypertarget{sourcepointer-7}{}\fi
\noindent\textit{Source reference.} Original reference 06WU: sources/stacks-chow/tags.txt, source lines 8828--8828. Cross-chapter standard prerequisite cited by the author inside the counted span; it is used as a pointer and its own proof is not part of this item.\par\medskip
\ifdefined\hypertarget\hypertarget{sourcepointer-8}{}\fi
\noindent\textit{Source reference.} Original lemma 0779: sources/stacks-chow/tags.txt, source lines 9194--9194. Cross-chapter standard prerequisite cited by the author inside the counted span; it is used as a pointer and its own proof is not part of this item.\par\medskip
\ifdefined\hypertarget\hypertarget{sourcepointer-9}{}\fi
\noindent\textit{Source reference.} Original lemma 07AR: sources/stacks-chow/tags.txt, source lines 9316--9316. Cross-chapter standard prerequisite cited by the author inside the counted span; it is used as a pointer and its own proof is not part of this item.\par\medskip
\begin{lemma}
\label{lemma-quasi-coherent}
Let $\mathcal{X}$ be an algebraic stack.
\begin{enumerate}
\item With $g$ as in Lemma \ref{lemma-lisse-etale}
for the lisse-\'etale site we have
\begin{enumerate}
\item the functors $g^{-1}$ and $g_!$ define mutually inverse functors
$$
\xymatrix{
\QCoh(\mathcal{O}_\mathcal{X}) \ar@<1ex>[r]^-{g^{-1}} &
\QCoh(\mathcal{X}_{lisse,\etale},
\mathcal{O}_{\mathcal{X}_{lisse,\etale}}) \ar@<1ex>[l]^-{g_!}
}
$$
\item if $\mathcal{F}$ is in $\textit{LQCoh}^{fbc}(\mathcal{O}_\mathcal{X})$
then $g^{-1}\mathcal{F}$ is in
$\QCoh(\mathcal{O}_{\mathcal{X}_{lisse,\etale}})$ and
\item $Q(\mathcal{F}) = g_!g^{-1}\mathcal{F}$ where $Q$ is as in
Lemma \ref{lemma-adjoint}.
\end{enumerate}
\item With $g$ as in Lemma \ref{lemma-lisse-etale}
for the flat-fppf site we have
\begin{enumerate}
\item the functors $g^{-1}$ and $g_!$ define mutually inverse functors
$$
\xymatrix{
\QCoh(\mathcal{O}_\mathcal{X}) \ar@<1ex>[r]^-{g^{-1}} &
\QCoh(\mathcal{X}_{flat,fppf},
\mathcal{O}_{\mathcal{X}_{flat,fppf}}) \ar@<1ex>[l]^-{g_!}
}
$$
\item if $\mathcal{F}$ is in $\textit{LQCoh}^{fbc}(\mathcal{O}_\mathcal{X})$
then $g^{-1}\mathcal{F}$ is in
$\QCoh(\mathcal{O}_{\mathcal{X}_{flat,fppf}})$
and
\item $Q(\mathcal{F}) = g_!g^{-1}\mathcal{F}$ where $Q$ is as in
Lemma \ref{lemma-adjoint}.
\end{enumerate}
\end{enumerate}
\end{lemma}

\begin{proof}
Pullback by any morphism of ringed topoi preserves categories of quasi-coherent
modules, see
Modules on Sites, Lemma \ref{sites-modules-lemma-local-pullback}.
Hence $g^{-1}$ preserves the categories of quasi-coherent modules;
here we use that
$\QCoh(\mathcal{O}_\mathcal{X}) =
\QCoh(\mathcal{X}_\etale, \mathcal{O}_\mathcal{X})$
by Sheaves on Stacks, Lemma
\ref{stacks-sheaves-lemma-characterize-quasi-coherent-bis}.
The same is true for $g_!$ by
Lemma \ref{lemma-shriek-quasi-coherent}.
We know that $\mathcal{H} \to g^{-1}g_!\mathcal{H}$ is an isomorphism by
Lemma \ref{lemma-lisse-etale}.
Conversely, if $\mathcal{F}$ is in $\QCoh(\mathcal{O}_\mathcal{X})$
then the map $g_!g^{-1}\mathcal{F} \to \mathcal{F}$ is a map of quasi-coherent
modules on $\mathcal{X}$ whose restriction to any scheme smooth over
$\mathcal{X}$ is an isomorphism. Then the discussion in
Sheaves on Stacks, Sections
\ref{stacks-sheaves-section-quasi-coherent-presentation} and
\ref{stacks-sheaves-section-quasi-coherent-algebraic-stacks}
(comparing with quasi-coherent modules on presentations)
shows it is an isomorphism. This proves (1)(a) and (2)(a).

\medskip\noindent
Let $\mathcal{F}$ be an object of
$\textit{LQCoh}^{fbc}(\mathcal{O}_\mathcal{X})$. By
Lemma \ref{lemma-adjoint-kernel-parasitic}
the kernel and cokernel of the map
$Q(\mathcal{F}) \to \mathcal{F}$ are parasitic. Hence by
Lemma \ref{lemma-parasitic-in-terms-flat-fppf}
and since $g^* = g^{-1}$ is exact, we conclude
$g^*Q(\mathcal{F}) \to g^*\mathcal{F}$ is an isomorphism. Thus
$g^*\mathcal{F}$ is quasi-coherent. This proves (1)(b) and (2)(b).
Finally, (1)(c) and (2)(c) follow because
$g_!g^*Q(\mathcal{F}) \to Q(\mathcal{F})$ is an isomorphism by
our arguments above.
\end{proof}

\end{document}
